\chapter{Conclusion}\label{cap:conclusion}

\section{Ce que je retiens du prototype}

J'ai examiné les contrôles des synthèses et des allocations proposées par le LLM.

Le backtest m'a fourni le socle de cette comparaison : l'équipondération des quatre ETF reste derrière SPY en rendement annualisé, après frais de transaction. Le contrôle sur 2021--2025 conserve ce classement. J'y vois un cas reproductible pour l'étude de la couche agentique, sans supériorité financière du multifactoriel.

Les synthèses DeepSeek satisfont les contrôles automatisés, mais trois altérations construites passent aussi le validateur : coût inexact, citation inexistante et recommandation contradictoire. Ces essais m'ont conduit à examiner ce que le validateur laisse passer, au-delà du taux d'acceptation. La fiabilité du texte et son utilité pour un analyste demanderaient également une lecture indépendante.

Dans l'extension d'allocation, DeepSeek termine derrière l'équipondération et l'inverse volatilité, après frais de transaction. Sur 56 propositions, 34 sont rejetées, surtout pour dépassement du turnover. Je retiens ici l'utilité du contrôle, qui bloque des propositions incorrectes. Le rendement obtenu dépend aussi de ces rejets et de la conservation des positions qu'ils déclenchent : il évalue l'ensemble du dispositif.

Ce travail m'a ainsi conduit à déplacer la question initiale : l'enjeu n'est pas seulement de savoir ce qu'un LLM peut produire pour un desk, mais surtout ce que le système est capable de vérifier avant d'utiliser cette production.

\section{Limites de portée}

Le PnL d'allocation reste rétrospectif : les indicateurs sont coupés dans le temps, mais le préentraînement peut contenir les événements étudiés. Une seule trajectoire a été collectée et les données ne sont pas disponibles dans leurs versions historiques point-in-time. Le replay et les 46 tests établissent la cohérence technique, sans démontrer la généralisation financière.

Je conserve également une réserve sur le bilan des coûts : le PnL exclut l'API, l'infrastructure et le travail humain. Les 66\,014 jetons connus, auxquels peut s'ajouter l'usage de la requête interrompue, documentent une consommation partielle. La facture, le coût total d'exploitation et les gains éventuels de temps restent à mesurer.

\section{Deux suites prioritaires}

Je renforcerais d'abord la fiabilité des sorties : tolérances adaptées aux champs, citations exactes et cohérence du texte pour les notes ; choix parmi des portefeuilles déjà admissibles pour l'allocation. Ce dernier dispositif serait comparé à une règle utilisant les mêmes indicateurs, selon un protocole défini avant observation des rendements puis suivi prospectivement.

Je comparerais ensuite les notes auprès d'analystes sur plusieurs dates, en faisant varier l'ordre de présentation. L'exactitude des réponses, la détection des alertes et le temps nécessaire seraient rapprochés des coûts API et de revue, pour évaluer la valeur pratique de cette intégration.
